% !TeX root = ../../../../../main.tex
% chapters/sections/chapter3/subsections/representative_household/euler_contingent.tex

\subsection{コンティンジェント債券に関するオイラー方程式}
\label{sec:chapter3_representative_household_euler_contingent}

\subsubsection{自国の自国コンティンジェント債券に関するオイラー方程式}
\label{sec:chapter3_representative_household_euler_contingent_home}

\paragraph{①自国の元となる方程式}
\label{sec:chapter3_representative_household_euler_contingent_home_original}
式
\eqref{form:chapter3_households_stage2_home_lambda_h_i_H_eq}
で示されたとおり、自国家計\(h\)の自国コンティンジェント債券\(d_{t+1}^{h\to H}(j')\)に関するFOCを、
\(j'\)について集計することで自国コンティンジェント債券\(d_{t+1}^{h\to H}(j')\)に関するオイラー方程式が
以下のように得られるのであった。
\begin{equation}
\label{form:chapter3_representative_household_euler_contingent_home_original_lambda_h_i_H_eq}
    \lambda_t^h=\beta_t^H(1+i_t^H)\operatorname{E}_t[\lambda_{t+1}^h]
\end{equation}

\paragraph{②自国の代表的家計の方程式}
\label{sec:chapter3_representative_household_euler_contingent_home_representative}
式
\eqref{form:chapter3_representative_household_euler_contingent_home_original_lambda_h_i_H_eq}
に対し
\eqref{form:chapter3_representative_household_lambda_home_representative_lambda_H_def}
により定義した自国代表的家計の所得の限界効用
\(\lambda_t^h=\lambda_t^H\)および
\(\lambda_{t+1}^h=\lambda_{t+1}^H\)
を代入すると自国代表的家計の自国コンティンジェント債券に関するオイラー方程式が得られる。
\begin{equation}
\label{form:chapter3_representative_household_euler_contingent_home_representative_lambda_H_i_H_eq}
    \lambda_t^H=\beta_t^H(1+i_t^H)\operatorname{E}_t[\lambda_{t+1}^H]
\end{equation}
さらに
\ref{chap:chapter5}
章において金融政策の効果を分析する準備として上式を前方展開しておく。
来期\(t+1\)においても同式が成立するため、それを代入し反復期待値の法則
\(\operatorname{E}_t[\operatorname{E}_{t+1}[\cdot]]=\operatorname{E}_t[\cdot]\)
を適用する操作を\(k\)期先まで繰り返すと以下が得られる。
\begin{equation}
\label{form:chapter3_representative_household_euler_contingent_home_representative_E_lambda_H_i_H_eq}
    \operatorname{E}_t[\lambda_{t+1}^H]
    =\operatorname{E}_t\left[\left(\prod_{k=1}^{k}\beta_{t+k}^H(1+i_{t+k}^H)\right)\lambda_{t+k+1}^H\right]
\end{equation}
本稿が使用するDynareによるシミュレーションにおいてはショックを受けた経済は
唯一の安定パス（サドルパス）をたどり新たな定常状態へ収束する。ただし実質変数は初期定常状態に収束する。
上式の両辺について極限\(\lim_{k\to\infty}\)をとると以下の
自国コンティンジェント債券に関する前方展開されたオイラー方程式が得られる。
\begin{equation}
\label{form:chapter3_representative_household_euler_contingent_home_representative_E_lambda_H_i_H_lim_eq}
    \begin{aligned}
        \operatorname{E}_t[\lambda_{t+1}^H]
        &=(\lim_{k\to\infty}\lambda_{t+k+1}^H)\operatorname{E}_t\left[\prod_{k=1}^{\infty}\beta_{t+k}^H(1+i_{t+k}^H)\right] \\
        &=(\lim_{k\to\infty}\frac{1}{p_{t+k+1}^{H\to W}c_{t+k+1}^{H\to W}})\operatorname{E}_t\left[\prod_{k=1}^{\infty}\beta_{t+k}^H(1+i_{t+k}^H)\right] \\
        &=(\frac{1}{(\lim_{k\to\infty}p_{t+k+1}^{H\to W})(\lim_{k\to\infty}c_{t+k+1}^{H\to W})})\operatorname{E}_t\left[\prod_{k=1}^{\infty}\beta_{t+k}^H(1+i_{t+k}^H)\right]
        &=(\frac{1}{(\lim_{k\to\infty}p_{t+k+1}^{H\to W})c_{ss}^{H\to W}})\operatorname{E}_t\left[\prod_{k=1}^{\infty}\beta_{t+k}^H(1+i_{t+k}^H)\right]
    \end{aligned}
\end{equation}

\subsubsection{外国の外国コンティンジェント債券に関するオイラー方程式}
\label{sec:chapter3_representative_household_euler_contingent_foreign}

\paragraph{①外国の元となる方程式}
\label{sec:chapter3_representative_household_euler_contingent_foreign_original}
式
\eqref{form:chapter3_households_stage2_foreign_lambda_f_slash_star_i_F_eq}
で示されたとおり、外国家計\(f\)の外国コンティンジェント債券\(d_{t+1}^{f\to F*}(j')\)に関する
FOCを\(j'\)について集計することで以下のオイラー方程式が得られるのであった。
\begin{equation}
\label{form:chapter3_representative_household_euler_contingent_foreign_original_lambda_f_slash_star_i_F_eq}
    \lambda_t^{f/*}=\beta_t^F(1+i_t^F)\operatorname{E}_t[\lambda_{t+1}^{f/*}]
\end{equation}

\paragraph{②外国の代表的家計の方程式}
\label{sec:chapter3_representative_household_euler_contingent_foreign_representative}
式
\eqref{form:chapter3_representative_household_euler_contingent_foreign_original_lambda_f_slash_star_i_F_eq}
に対し
\eqref{form:chapter3_representative_household_lambda_foreign_representative_lambda_F_slash_star_def}
により定義した外国代表的家計の所得の限界効用
\(\lambda_t^{f/*}=\lambda_t^{F/*}\)および
\(\lambda_{t+1}^{f/*}=\lambda_{t+1}^{F/*}\)
を代入すると外国代表的家計の外国コンティンジェント債券に関するオイラー方程式が得られる。
\begin{equation}
\label{form:chapter3_representative_household_euler_contingent_foreign_representative_lambda_F_slash_star_i_F_eq}
    \lambda_t^{F/*}=\beta_t^F(1+i_t^F)\operatorname{E}_t[\lambda_{t+1}^{F/*}]
\end{equation}
さらに
\ref{chap:chapter5}
章において金融政策の効果を分析する準備として上式を前方展開しておく。
来期\(t+1\)においても同式が成立するため、それを代入し反復期待値の法則
\(\operatorname{E}_t[\operatorname{E}_{t+1}[\cdot]]=\operatorname{E}_t[\cdot]\)
を適用する操作を\(k\)期先まで繰り返すと以下が得られる。
\begin{equation}
\label{form:chapter3_representative_household_euler_contingent_foreign_representative_E_lambda_F_slash_star_i_F_eq}
    \operatorname{E}_t[\lambda_{t+1}^{F/*}]
    =\operatorname{E}_t\left[\left(\prod_{k=1}^{k}\beta_{t+k}^F(1+i_{t+k}^F)\right)\lambda_{t+k+1}^{F/*}\right]
\end{equation}
本稿が使用するDynareによるシミュレーションにおいてはショックを受けた経済は
唯一の安定パス（サドルパス）をたどり新たな定常状態へ収束する。ただし実質変数は初期定常状態に収束する。
上式の両辺について極限\(\lim_{k\to\infty}\)をとると以下の
外国コンティンジェント債券に関する前方展開されたオイラー方程式が得られる。
\begin{equation}
\label{form:chapter3_representative_household_euler_contingent_foreign_representative_E_lambda_F_slash_star_i_F_lim_eq}
    \begin{aligned}
        \operatorname{E}_t[\lambda_{t+1}^{F/*}]
        &=(\lim_{k\to\infty}\lambda_{t+k+1}^{F/*})\operatorname{E}_t\left[\prod_{k=1}^{\infty}\beta_{t+k}^F(1+i_{t+k}^F)\right] \\
        &=(\lim_{k\to\infty}\frac{1}{p_{t+k+1}^{F\to W*}c_{t+k+1}^{F\to W}})\operatorname{E}_t\left[\prod_{k=1}^{\infty}\beta_{t+k}^F(1+i_{t+k}^F)\right] \\
        &=(\frac{1}{(\lim_{k\to\infty}p_{t+k+1}^{F\to W*})(\lim_{k\to\infty}c_{t+k+1}^{F\to W})})\operatorname{E}_t\left[\prod_{k=1}^{\infty}\beta_{t+k}^F(1+i_{t+k}^F)\right] \\
        &=(\frac{1}{(\lim_{k\to\infty}p_{t+k+1}^{F\to W*})c_{ss}^{F\to W}})\operatorname{E}_t\left[\prod_{k=1}^{\infty}\beta_{t+k}^F(1+i_{t+k}^F)\right]
    \end{aligned}
\end{equation}